% Source: https://artofproblemsolving.com/community/c3h3422533_one_nice
% Stable identity: aops-topic-3422533-q01
Occurrences wholly inside one number comprise the number $2022$, the nine numbers $d2022$ with $1\le d\le9$, and the ten numbers $20220$ through $20229$: twenty in total. For occurrences crossing one boundary after two digits, the preceding number ends in $20$ and its successor begins with $22$. These preceding numbers are $220,2220$, and $22020,22120,\ldots,22920$, giving twelve. For a boundary after three digits, the preceding number ends in $202$ and its successor begins with $2$: these are $202,2202$, and $20202,21202,\ldots,29202$, another twelve. A boundary after one digit would require a successor beginning with $022$, impossible. An occurrence spanning two boundaries must contain an entire one- or two-digit integer in its interior. One-digit cases lie near the start and direct inspection excludes them; the only possible two-digit interior strings are $02$ and $20$. The former is not a representation, and the predecessor of $20$ ends in $9$, not $2$. Thus the total is $\boxed{44}$.
